\documentclass[12pt, a4paper]{article}
\usepackage[utf8]{inputenc}
\usepackage[T2A]{fontenc}
\usepackage{indentfirst, setspace}
\usepackage{tabularx, multirow}
\usepackage[normalem]{ulem}
\usepackage[style=russian]{csquotes}
\usepackage[english,russian]{babel}
\usepackage{hyperref}
\usepackage{ragged2e}
\usepackage{caption}
\usepackage{wrapfig}
\usepackage{amsmath}
\usepackage{tikz}
\makeatletter
\def\@biblabel#1{#1. }
\makeatother
\captionsetup{labelsep=endash}
\usepackage{listings}
\linespread{1.3}
\lstset{
  language=C++, % Язык программирования (можно выбрать другой)
  basicstyle=\ttfamily\small, % Основной стиль шрифта
  keywordstyle=\color{blue}, % Стиль ключевых слов
  breaklines=true,
  commentstyle=\color{green}, % Стиль комментариев
  stringstyle=\color{red}, % Стиль строк
  extendedchars=\true,
  showstringspaces=false,
}
\usepackage[left=20mm, top=15mm, right=15mm, bottom=20mm]{geometry}

\begin{document}
\begin{titlepage}
     \fontsize{12}{12}\selectfont

  {\centering

   \begin{bf}

    \begin{wrapfigure}{l}{19mm}
        
\includegraphics[width=22mm]{emblem.jpeg}
    \end{wrapfigure}


    \noindent Министерство науки и высшего образования Российской Федерации

    \noindent Федеральное государственное бюджетное образовательное учреждение высшего образования

    \noindent \enquote{Московский государственный технический унверситет

     \noindent имени Н.Э. Баумана

     \noindent (национальный исследовательский университет)}

    \noindent (МГТУ им. Н.Э. Баумана)

   \end{bf}
  }

  \vspace{0.4cm}

  {\setstretch{0.1}
   \noindent\rule{\textwidth}{1mm}
   \noindent\rule{\textwidth}{0.5mm}

  }

  \fontsize{14}{21}\selectfont

  \noindent\begin{tabularx}{\textwidth}{l >{\centering\arraybackslash}X}
   ФАКУЛЬТЕТ & \flqq Фундаментальные Науки\frqq \\ \cline{2-2}

   КАФЕДРА & ФН-12 \flqq Математическое моделирование\frqq \\ \cline{2-2}
  \end{tabularx}


  \vspace{1cm}


  \begin{center}
   \begin{bf}

    \fontsize{24}{36}\selectfont
    ОТЧЕТ

    \fontsize{20}{30}\selectfont
    ПО ЛАБОРАТОРНОЙ РАБОТЕ НА ТЕМУ:

    Обход графовых структур

   \end{bf}
  \end{center}

  \fontsize{14}{21}\selectfont
  \vspace{5cm}


  \noindent\begin{tabularx}{\textwidth}{ X >{\centering}p{4cm} p{1cm} c }
   Студент: & & & Подобедов В. В. \\ \cline{2-2} \cline{4-4}
   & \fontsize{10}{15}\selectfont дата, подпись & & \fontsize{10}{15}\selectfont Ф.И.О. \\
   Преподаватель: & & & Волкова Л. Л.\\ \cline{2-2} \cline{4-4}
   & \fontsize{10}{15}\selectfont дата, подпись & & \fontsize{10}{15}\selectfont Ф.И.О.
  \end{tabularx}

  \vspace{\fill}

  \begin{center}
   \it{Москва}, 2023
  \end{center}

  \thispagestyle{empty}
\end{titlepage}

\renewcommand{\contentsname}{Содержание}
\setcounter{page}{2}
\tableofcontents

\newpage
\begin{center}
   \section*{Введение} 
\addcontentsline{toc}{section}{Введение}
\end{center}
\justifying
\textbf{Цель лабораторной работы}: реализовать алгоритмы обхода графовых структур.

Для достижения поставленной цели требуется решить следующие \textbf{задачи}.
\begin{enumerate}
\item Описать граф и бинарное дерево.
\item Реализовать графовую структуру.
\item Реализовать возможность пользователю самостоятельно заполнять графовую структуру.
\item Реализовать обход графовой структуры в ширину и в глубину.
\item Выполнить тестирование реализации разработанного алгоритма.
\end{enumerate}

Согласно варианту, требуется разработать бинарное дерево.

\newpage
\begin{center}
   \section{Аналитическая часть} 
\end{center}

\textbf{Граф} --- это математическая абстракция, представляющая собой набор вершин и рёбер между этими вершинами.

\textbf{Граф с петлями} --- это граф, в котором рёбра могут соединять вершины сами с собой, создавая так называемые "петли". Петля представляет собой ребро, начало и конец которого находятся в одной и той же вершине.

\textbf{Граф без петель} --- это граф, в котором отсутствуют рёбра, соединяющие вершину с самой собой.

\textbf{Бинарное дерево} --- это иерархическая структура данных в виде дерева, в которой каждый узел может иметь не более двух дочерних узлов: левый и правый. У дерева есть две главные характеристики:
\begin{enumerate}
    \item Корень --- верхний узел дерева, от которого начинаются все другие узлы. У бинарного дерева может быть только один корень.
    \item Листья --- узлы без дочерних узлов называются листьями. Листья находятся на самом нижнем уровне дерева.
\end{enumerate}

\newpage
\begin{center}
    \section{Конструкторская часть}
\end{center}
\subsection{Обход бинарного дерева в ширину}
Обход бинарного дерева в ширину происходит по уровням дерева: сначала посещаются все узлы на одном уровне, затем переходят к узлам следующего уровня. Для этого используется структура данных очередь (queue), которая помогает сохранять порядок обработки узлов.

Алгоритм обхода бинарного дерева в ширину:
\begin{enumerate}
    \item Начать с корневого узла дерева и поместить его в очередь.
    \item Пока очередь не пуста:
    \begin{itemize}
        \item[---] Извлечь узел из начала очереди.
        \item[---] Посетить этот узел и напечатать его значение.
        \item[---] Если у узла есть левый потомок, поместить его в конец очереди.
        \item[---] Если у узла есть правый потомок, поместить его в конец очереди.
    \end{itemize}
    \item Повторить шаг 2, пока очередь не опустеет.
\end{enumerate}

\subsection{Обход бинарного дерева в глубину}
Обход бинарного дерева в глубину начинается с корневого узла и идет вглубь структуры дерева, прежде чем вернуться обрабатывать более поздние узлы. Для этого используется структура данных стек (stack), которая помогает сохранять порядок обработки узлов.

Алгоритм обхода бинарного дерева в ширину:
\begin{enumerate}
    \item Начать с корневого узла дерева и поместить его в стек.
    \item Пока стек не пуст:
    \begin{itemize}
        \item[---] Извлечь узел из начала стека.
        \item[---] Посетить этот узел и напечатать его значение.
        \item[---] Если у узла есть правый потомок, поместить его в конец стека.
        \item[---] Если у узла есть левый потомок, поместить его в конец стека.
    \end{itemize}
    \item Повторить шаг 2, пока стек не опустеет.
\end{enumerate}

\subsection{Написание бинарного дерева}
С помощью шаблонных классов ``BinaryTree'', который описывает саму структуру бинарного дерева, и ``BinaryTreeNode'', который описывает узел бинарного дерева, реализовать бинарное дерево на ссылках. При этом нужно учесть, что левый или правый потомок могут отсутствовать.  Требуется разработать следующие методы.
\begin{itemize}
    \item[---] Рекурсивный метод ``addElement'' --- метод, который добавляет элемент в бинарное дерево.
    \item[---] Рекурсивный метод ``BreadthTree'' --- метод обхода бинарного дерева в ширину.
    \item[---] Рекурсивный метод ``DepthTree'' --- метод обхода бинарного дерева в глубину.
    \item[---] Метод ``addElementToTree'' --- метод, который обрабатывает введённый пользователем элемент и вызывает рекурсивный метод ``addElement''.
    \item[---] Метод ``printBreadthTree'' --- метод, который выводит в консоль результат обхода дерева в ширину.
    \item[---] Метод ``printDepthTree'' --- метод, который выводит в консоль результат обхода дерева в глубину.
\end{itemize}

К программе предъявлены следующие требования. Пользователь самостоятельно вводит данные в бинарное дерево, после чего выводится результат обхода в ширину и глубину пользовательского дерева.

\newpage
\begin{center}
    \section{Технологическая часть}
\end{center}

Лабораторная работа была реализована на языке программирования C++. Кроме стандартной библиотеки <iostream>, которая предоставляет возможность пользоваься средствами ввода-вывода для стандартной консоли, были подключены библиотеки <queue> и <stack>. Они представляют собой контейнеры, реализующие структуры данных ``стек'' и ``очередь'' соответственно, предоставляя удобные интерфейсы для хранения и управления данными в виде стека или очереди. Также была подключена библиотека <string> для пользования строками. Был также написан стек на ссылках.

Далее приведен листинг реализации стека (листинг~\ref{lst:label1}).
\begin{lstlisting}[caption={Исходный код программы для реализации стека},label={lst:label1}]
#include <iostream>
#include <string>
#include <queue>
#include <stack>

using namespace std;

template <typename T>
class BinaryTreeNode {
public:
	T data;
	BinaryTreeNode<T>* left;
	BinaryTreeNode<T>* right;

	BinaryTreeNode(T value) : data(value), left(nullptr), right(nullptr) {}
};

template <typename T>
class BinaryTree {
private:
	BinaryTreeNode<T>* root;

	void addElement(BinaryTreeNode<T>* node) {
		string choice;

		cout << "Добавить значение в левый узел для узла " << node->data << "? (yes/no): ";
		cin >> choice;

		if (choice == "yes") {
			T value;
			cout << "Введите значение для левого узла: ";
			cin >> value;
			node->left = new BinaryTreeNode<T>(value);
			addElement(node->left);
		}
		else if (choice == "no") {
			node->left = nullptr;
		}

		cout << "Добавить значение в правый узел для узла " << node->data << "? (yes/no): ";
		cin >> choice;

		if (choice == "yes") {
			T value;
			cout << "Введите значение для правого узла: ";
			cin >> value;
			node->right = new BinaryTreeNode<T>(value);
			addElement(node->right);
		}
		else if (choice == "no") {
			node->right = nullptr;
		}
	}

	void BreadthTree(BinaryTreeNode<T>* node) {
		if (node == nullptr) {
			return;
		}

		queue<BinaryTreeNode<T>*> nodesQueue;
		nodesQueue.push(node);

		while (!nodesQueue.empty()) {
			BinaryTreeNode<T>* current = nodesQueue.front();
			nodesQueue.pop();
			cout << current->data << " ";

			if (current->left != nullptr) {
				nodesQueue.push(current->left);
			}
			if (current->right != nullptr) {
				nodesQueue.push(current->right);
			}
		}
	}

	void DepthTree(BinaryTreeNode<T>* node) {
		if (node == nullptr) {
			return;
		}

		stack<BinaryTreeNode<T>*> nodesStack;
		nodesStack.push(node);

		while (!nodesStack.empty()) {
			BinaryTreeNode<T>* current = nodesStack.top();
			nodesStack.pop();
			std::cout << current->data << " ";

			if (current->right != nullptr) {
				nodesStack.push(current->right);
			}
			if (current->left != nullptr) {
				nodesStack.push(current->left);
			}
		}
	}

public:
	BinaryTree() : root(nullptr) {}

	void addElementToTree() {
		T value;
		cout << "Введите значение для корня дерева: ";
		cin >> value;
		root = new BinaryTreeNode<T>(value);
		addElement(root);
	}

	void printBreadthTree() {
		cout << "Обход дерева в ширину: ";
		BreadthTree(root);
		cout << endl;
	}

	void printDepthFirstTree() {
		cout << "Обход дерева в глубину: ";
		DepthTree(root);
		cout << endl;
	}
};

int main() {
	setlocale(0, "");
	BinaryTree<int> binaryTree;
	binaryTree.addElementToTree();
	cout << endl;
	binaryTree.printBreadthTree();
	cout << endl;
	binaryTree.printDepthFirstTree();
	return 0;
}
\end{lstlisting}

\newpage
\begin{center}
    \section{Исследовательская часть}
\end{center}

Далее на рисунках представлен результат работы программы и входной файл с инфиксным выражением.

Результат работы программы, где пользователь вводит бинарное дерево, и программа выводит обход пользовательского дерева в ширину и глубину. Результат приведён на рис.~\ref{img:grap1}. 
\begin{figure}[h!]
    \center{\includegraphics{test1.PNG}}
    \caption{Пример №1}
    \label{img:grap1}
\end{figure}

\newpage
Результат работы программы, где пользователь вводит бинарное дерево, состоящее только из левых потомков, и программа выводит обход пользовательского дерева в ширину и глубину. Результат приведён на рис.~\ref{img:grap2}. 
\begin{figure}[h!]
    \center{\includegraphics{test2.PNG}}
    \caption{Пример №2}
    \label{img:grap2}
\end{figure}

\newpage
\begin{center}
    \section*{Заключение}
    \addcontentsline{toc}{section}{Заключение}
\end{center}

Цель лабораторной работы достигнута: реализован алгоритмы обхода бинарного дерева.

В результате выполнения лабораторной работы были выполнены все задачи.
\begin{enumerate}
\item Описано бинарное дерево.
\item Программно реализовано бинарное дерево.
\item Реализована возможность пользователю самостоятельно заполнять бинарное дерево.
\item Реализован обход бинарного дерева в ширину и в глубину.
\item Выполнено тестирование реализации разработанного алгоритма.
\end{enumerate}

\newpage
\begin{center}
\begin{thebibliography}{}
\addcontentsline{toc}{section}{Список используемых источников}
\bibitem{book} Рафгарден Тим. Совершенный алгоритм. Графовые алгоритмы и структуры данных. --- СПб.: Питер, 2019. --- 256 с.: ил. --- (Серия ``Библиотека программиста'').
\end{thebibliography}
\end{center}

\end{document}
